% !TEX program = xelatex
% =====================================================================
%  圆锥曲线的三维矢量方法 —— 导言区
%  编译方式：xelatex main.tex  （连续编译两次以生成目录与交叉引用）
% =====================================================================

\usepackage{amsmath,amssymb,amsthm}
\usepackage{geometry}
\usepackage{mathrsfs}
\usepackage{array}
\usepackage{booktabs}
\usepackage{enumitem}
\usepackage{tikz}
\usetikzlibrary{calc,intersections}
\usepackage[hidelinks,bookmarksnumbered]{hyperref}

\geometry{a4paper,left=2.6cm,right=2.6cm,top=2.8cm,bottom=2.8cm}

% ---------- 定理类环境 ----------
\theoremstyle{definition}
\newtheorem{definition}{定义}[chapter]
\newtheorem{example}[definition]{例}
\newtheorem{exercise}[definition]{习题}
\newtheorem{remark}[definition]{注}
\theoremstyle{plain}
\newtheorem{theorem}[definition]{定理}
\newtheorem{proposition}[definition]{命题}
\newtheorem{lemma}[definition]{引理}
\newtheorem{corollary}[definition]{推论}

\renewcommand{\proofname}{\bfseries 证明}
\renewcommand{\contentsname}{目\quad 录}

% ---------- 常用记号 ----------
\newcommand{\R}{\mathbb{R}}
\newcommand{\dd}{\mathrm{d}}
\newcommand{\ii}{\mathrm{i}}
\newcommand{\tr}{\operatorname{tr}}
\newcommand{\diag}{\operatorname{diag}}
\newcommand{\rank}{\operatorname{rank}}
\newcommand{\sgn}{\operatorname{sgn}}
\DeclareMathOperator{\Ker}{Ker}
\DeclareMathOperator{\Ima}{Im}

% 度规无关记号：混合积
\newcommand{\mix}[3]{[#1\,#2\,#3]}
% 三阶反对称张量
\newcommand{\eps}{\epsilon}
% 对称积 (\cdot\cdot)
\newcommand{\sym}[2]{(#1#2)}
% 无穷远线矢
\newcommand{\linf}{l_{\infty}}
\newcommand{\linfsub}[1]{l_{\infty #1}}

\allowdisplaybreaks

% 中文段落断行容忍度：避免个别行溢出页面右边界
\setlength{\emergencystretch}{2.5em}
\tolerance=2000
\hbadness=3000
